% !PRINTED 66
\weberpage{66}
%<<CONT
3. \gesperrt{Between at most $m+1$ related $T$-functions of order $m$, $T, T_1, \ldots, T_m$, there exists an identical equation of the form}
\[
LT + L_1T_1 + \cdots + L_mT_m = 0.
\]
For if this relation does not already hold for $L=0$, then among the constants available in $L_1, L_2, \ldots, L_m$ one can first determine $\lambda_1, \lambda_2 \ldots, \lambda_m, \mu_1, \mu_2 \ldots, \mu_m$ in such a way that all the products $L_1T_1, L_2T_2, \ldots, L_mT_m$ receive the same periodicity factors, and then the $C_1, C_2, \ldots, C_m$ in such a way that $m-1$ of the zero points of the function $L_1T_1+L_2T_2+\cdots+L_mT_m$, and consequently all of them, coincide with the zero points of the function $T$. From this, however, it follows that the ratio of $L_1T_1+L_2T_2+\cdots+L_mT_m$ to $T$ is equal to a $T^0$-function $-L$, which is what was to be proved.

4. \gesperrt{The quotients of related $T$-functions can be transformed into doubly periodic functions by the addition of a factor $L$.}

5. \gesperrt{If we assume the existence of a $T$-function of the first order $t(u)$, then every $T$-function of order $m$ can be derived from it.}

For let $\gamma$ be the character of $t(u)$, and let $\alpha_1, \alpha_2, \ldots, \alpha_m$ be arbitrarily given values. Then the function
\[
t(u-\alpha_i+\gamma) = t_i(u), \qquad i = 1, 2 \ldots m
\]
vanishes at the point $\alpha_i$ and at all points congruent to $\alpha_i$, but at no other.

The product
\[
T(u) = t_1(u)t_2(u) \ldots t_m(u)
\]
is therefore a $T$-function of order $m$ with the arbitrarily given zero points $\alpha_1, \alpha_2, \ldots, \alpha_m$.

6. \gesperrt{From this it can easily be proved that every doubly periodic function which has only a finite number of points of discontinuity in a period parallelogram can be represented as the quotient of two $T$-functions.}

For let $\varphi(u)$ be a doubly periodic function with the periods $\omega_1, \omega_2$ and the points of discontinuity $\alpha_1, \alpha_2, \ldots, \alpha_m$, which may also in part coincide in points of discontinuity of higher order.
